\documentclass[11pt]{article}

\usepackage[margin=1in]{geometry}
\usepackage{amsmath, amssymb, amsthm}
\usepackage{bm}
\usepackage{hyperref}
\usepackage{enumitem}
\usepackage{authblk}

\title{Prime-Structured Tensor-Network Autoencoder:\\
Architecture, Regularization, and Prior Art Defensive Publication}

\author{Citizen Gardens \\ Prime Materia Commons}
\date{May 3, 2026}

\begin{document}
\maketitle

\begin{abstract}
This report presents a mathematically consistent and computationally implementable framework for \emph{prime-structured tensor-network autoencoders}. The framework begins from a classical tensor-network (TN) autoencoder baseline, then introduces (i) prime-structured design of allowed bond dimensions, and (ii) differentiable regularizers that bias effective entanglement ranks toward this prime-structured lattice of dimensions. This document also delineates the boundary between discrete architectural choices and continuous parameter optimization, clarifies the role of prime indexing, and proposes both weak and strong forms of prime-aware regularization. The goal is to establish prior art and provide a reference point for future developments in multiplicity-aware tensor networks, including their potential quantum or hybrid implementations.
\end{abstract}

\tableofcontents
\newpage
\section{Executive Summary}

This document consolidates and refines the development of a novel architectural and regularization framework for tensor-network autoencoders (TN-AEs), with a focus on \emph{prime-structured bond dimensions}. The core object is NN-AUTOTENSCODER, a classical tensor-network autoencoder whose design and training explicitly exploit the prime factorization of bond dimensions.

\subsection{Key Ideas}

\begin{itemize}[leftmargin=1.5em]
    \item \textbf{Classical TN Autoencoder Baseline.} The Autoencoder is first defined as a purely classical TN autoencoder (MPS/TTN-based), with fixed topology, differentiable tensor entries, and a standard reconstruction loss. All gradients are obtained via automatic differentiation or backpropagation through tensor contractions.

    \item \textbf{Prime-Structured Architectural Design.} Each bond in the encoder and decoder TN is constrained to take values from a finite set of allowed bond dimensions, each with a prescribed prime factorization. The \emph{prime structure} thus enters the model as a design constraint on representable entanglement ranks.

    \item \textbf{Effective Rank via Singular Values.} For each intermediate state in the encoder, the effective bond dimension (rank) is extracted by canonicalizing the corresponding TN and inspecting singular values across each bond. This is standard in TN practice and ties the framework directly to entanglement structure.

    \item \textbf{Differentiable Rank Surrogates and Regularization.} A smooth surrogate for effective rank is constructed from the singular values using sigmoid-type functions. A prime-aware regularizer then penalizes deviations of this surrogate rank from the allowed discrete bond dimensions. A softmin variant ensures differentiability.

    \item \textbf{Weak vs.\ Strong Prime Regularization.} In the weak form, prime structure is encoded solely through the set of allowed integers. In the strong form, approximate prime exponent vectors are extracted via a log-based regression, and their deviation from target exponent patterns is explicitly penalized.

    \item \textbf{Architectural Multiplicity Schedules.} Rather than dynamically splitting or fusing bonds, multiplicity moves are interpreted as design patterns in the allowed bond dimension sets across layers, encoding capacity shrinkage (compression) or expansion.

    \item \textbf{Prime-Canonical Model Selection.} A complexity functional based on prime exponents of maximal allowed bond dimensions is proposed as a criterion for architecture selection, preferring simpler prime-structured entanglement budgets subject to reconstruction performance.

    \item \textbf{Prior-Art Intent.} This document is intended both as a substantive technical report and as a defensive publication establishing prior art around prime-indexed bond design, prime-aware rank regularization, and multiplicity-style architecture selection in TN autoencoders.
\end{itemize}

\subsection{What is New Compared to Standard TN-AEs}

While tensor-network autoencoders and rank-regularized models are known in the literature, the specific combination of:

\begin{itemize}[leftmargin=1.5em]
    \item constraining bond dimensions to prime-structured sets;
    \item using differentiable surrogates of effective rank aligned to those sets; and
    \item optionally defining and regularizing approximate prime exponent vectors,
\end{itemize}
appears to be novel as a unified framework. This prior art claim applies particularly to settings where these design choices are used explicitly to interpret and bias entanglement structure in terms of prime-indexed multiplicity channels.

\section{Baseline: Classical Tensor-Network Autoencoder}

We first define a classical tensor-network autoencoder without prime structure.

\subsection{Setting and Notation}

Let
\[
\mathcal{H} \;=\; (\mathbb{C}^d)^{\otimes N}
\]
be the Hilbert space of an $N$-site system (e.g., $d=2$ for qubits or binary features). A data set is a finite collection
\[
\mathcal{D} = \{\lvert x^{(n)} \rangle \in \mathcal{H}\}_{n=1}^M.
\]

We assume that elements of $\mathcal{D}$ are either directly represented as tensor networks (e.g.\ MPS) or can be mapped to such representations.

\subsection{Encoder and Decoder as Tensor Networks}

A \emph{tensor-network autoencoder} consists of:

\begin{itemize}[leftmargin=1.5em]
    \item An \textbf{encoder} $E_\theta : \mathcal{H} \to \mathcal{H}_{\mathrm{lat}}$, implemented as a fixed-topology TN (MPS, TTN, MERA-like, etc.), with trainable tensor entries collectively denoted by parameters $\theta$.

    \item A \textbf{decoder} $D_\theta : \mathcal{H}_{\mathrm{lat}} \to \mathcal{H}$, also a fixed-topology TN, either sharing parameters with $E_\theta$ or with its own parameter set.
\end{itemize}

For each input $\lvert x \rangle$, the encoded and reconstructed states are
\[
\lvert z_\theta(x) \rangle = E_\theta \lvert x \rangle, 
\qquad
\lvert \hat{x}_\theta \rangle = D_\theta \lvert z_\theta(x) \rangle.
\]

\subsection{Reconstruction Loss and Differentiability}

For a reconstruction objective, we may choose a fidelity-like loss:
\[
\mathcal{L}_{\mathrm{rec}}(\theta)
= \frac{1}{|\mathcal{D}|} \sum_{\lvert x \rangle\in \mathcal{D}}
\left(
1 - \frac{|\langle x \mid \hat{x}_\theta \rangle|^2}{\|x\|^2 \|\hat{x}_\theta\|^2}
\right).
\]
Other choices (e.g., $\|\lvert x \rangle - \lvert \hat{x}_\theta \rangle\|^2$) are also possible.

The tensor entries depend smoothly on $\theta$; TN contraction is a differentiable transformation when implemented in standard autodiff frameworks. Thus $\mathcal{L}_{\mathrm{rec}}(\theta)$ is differentiable almost everywhere, and gradients $\nabla_\theta \mathcal{L}_{\mathrm{rec}}$ can be obtained via automatic differentiation.

\section{Prime-Structured Architectural Design}

We now introduce the central idea: \emph{prime-structured bond dimensions as an architectural constraint}.

\subsection{Bond Set and Allowed Dimensions}

Let $\mathcal{B}$ denote the set of virtual bonds (internal indices) in the encoder TN graph. For each bond $b\in \mathcal{B}$, we specify a finite set of allowed bond dimensions
\[
\mathcal{D}_b \;\subset\; \mathbb{N}.
\]

Each allowed dimension $D\in\mathcal{D}_b$ is required to admit a prime factorization over a prescribed finite set of primes $\mathcal{P}_b = \{p_{b,1},\dots,p_{b,M_b}\}$:
\[
D \;=\; \prod_{j=1}^{M_b} p_{b,j}^{k_{b,j}(D)},
\quad 
k_{b,j}(D)\in\mathbb{N}_0.
\]

We define the \emph{prime profile} associated to $D$ by
\[
\mathbf{k}_b(D) \;:=\; (k_{b,1}(D),\dots,k_{b,M_b}(D)) \in \mathbb{N}_0^{M_b}.
\]

The architecture class of TN encoders (and similarly decoders) is then:
\[
\mathcal{A} = \{ E_\theta:
\text{TN with fixed graph, and bond dimensions } D_b \in\mathcal{D}_b \text{ for all } b\in\mathcal{B} \}.
\]

In other words, the TN topology is fixed, and bond dimensions are constrained to chosen prime-structured discrete sets, not dynamically created or destroyed during training.

\subsection{Multiplicity Schedules Across Layers}

If the encoder is layered, $E_\theta = E_\theta^{(L)}\circ\cdots\circ E_\theta^{(1)}$, we may define layer-dependent sets $\mathcal{D}_b^{(\ell)}$ for bonds at layer~$\ell$:
\[
D_b^{(\ell)} \in \mathcal{D}_b^{(\ell)}.
\]

Multiplicity moves (e.g.\ ``fuse'' or ``split'') are then \emph{not} operations enacted by tensors on-the-fly, but rather design patterns such as:
\begin{itemize}[leftmargin=1.5em]
    \item \textbf{Fuse schedule:} $\mathcal{D}^{(\ell+1)}_b \subseteq \mathcal{D}^{(\ell)}_b$ --- entanglement capacity shrinks with depth.
    \item \textbf{Split schedule:} $\mathcal{D}^{(\ell+1)}_b \supseteq \mathcal{D}^{(\ell)}_b$ --- capacity expands (e.g.\ toward decoder side).
\end{itemize}

These patterns can be chosen to respect divisibility or prime-exponent relationships. For instance, in a compression region, we may require that each $D' \in \mathcal{D}_b^{(\ell+1)}$ divides some $D \in \mathcal{D}_b^{(\ell)}$, encoding a notion of multiplicity reduction.
\newpage
\section{Extracting Effective Ranks via Canonicalization}

To tie the prime-structured architecture to actual entanglement usage, we inspect the effective ranks of intermediate representations in the encoder.

\subsection{Hierarchical Encoders and Layers}

We assume either:
\begin{itemize}[leftmargin=1.5em]
    \item A hierarchical encoder (TTN/MERA-like), where each ``layer'' corresponds to a level in the tree or causal cone, and internal bonds at that level are naturally defined; or

    \item An MPS/MPO-based encoder realized as a sequence of MPO layers, such that after each layer we canonically represent the resulting state as an MPS with well-defined bonds.
\end{itemize}

In either case, after applying a layer $E^{(\ell)}_\theta$, we perform a canonicalization step on the resulting TN state to obtain singular values across each bond.

\subsection{Canonical Form and Singular Values}

For concreteness, consider MPS encoders. A state $\lvert \psi \rangle$ represented as an MPS can be brought into mixed canonical form. For each internal bond $b$, there exists a Schmidt decomposition
\[
\lvert \psi \rangle = \sum_{j=1}^{R_b} \sigma_{b,j} \, \lvert L_{b,j} \rangle \otimes \lvert R_{b,j} \rangle,
\]
where $\sigma_{b,j}\ge 0$ are singular values, $R_b$ is the rank, and $\{\lvert L_{b,j} \rangle\}$, $\{\lvert R_{b,j} \rangle\}$ are orthonormal. Libraries such as ITensor, TeNPy, and TensorNetwork provide standard routines for this canonicalization and SVD extraction.

Thus, for a given input $\lvert x \rangle$ and layer $\ell$, we obtain for each bond $b$ the set $\{\sigma_{b,j}^{(\ell)}(\theta,x)\}_{j}$.

\subsection{Effective Rank and Rank Surrogates}

Define the \emph{effective rank} at threshold $\varepsilon>0$:
\[
r_b^{(\ell)}(\theta,x)
:= \#\{j: \sigma_{b,j}^{(\ell)}(\theta,x) > \varepsilon\}.
\]
This is a discrete integer. To make it differentiable, we introduce a smooth surrogate
\[
\tilde{r}_b^{(\ell)}(\theta,x)
:= \sum_{j} s\left( \frac{\sigma_{b,j}^{(\ell)}(\theta,x)}{\sigma_{b,\max}^{(\ell)}(\theta,x)}\right),
\]
where $\sigma_{b,\max}^{(\ell)} = \max_j \sigma_{b,j}^{(\ell)}$ and
\[
s(u) = \frac{1}{1 + e^{-\alpha (u - \tau)}},\quad u\in[0,1],
\]
with $\alpha>0$ controlling steepness and $\tau\in(0,1)$ a relative threshold. This $\tilde{r}_b^{(\ell)}$ is a differentiable approximation to the number of ``significant'' singular values.

\section{Prime-Aware Regularization: Weak Form}

We now define the \emph{weak} prime-aware regularizer, which uses only the discrete allowed bond dimensions and the rank surrogate.

\subsection{Soft Distance to Prime-Structured Dimensions}

For each bond $b$ with allowed dimensions $\mathcal{D}_b = \{D_{b,1},\dots,D_{b,m_b}\}$, define a soft distance
\begin{align*}
w_{b,i}^{(\ell)}(\theta,x)
&:= \frac{\exp(-\beta \lvert \tilde{r}_b^{(\ell)}(\theta,x) - D_{b,i}\rvert)}
{\sum_{k=1}^{m_b} \exp(-\beta \lvert \tilde{r}_b^{(\ell)}(\theta,x) - D_{b,k}\rvert)},\\
\mathrm{dist}_b^{(\ell)}(\theta,x)
&:= \sum_{i=1}^{m_b}
w_{b,i}^{(\ell)}(\theta,x)\,\lvert \tilde{r}_b^{(\ell)}(\theta,x) - D_{b,i}\rvert,
\end{align*}
with $\beta>0$ a sharpness parameter. As $\beta\to\infty$, $\mathrm{dist}_b^{(\ell)}$ approaches $\min_{D\in\mathcal{D}_b} \lvert \tilde{r}_b^{(\ell)} - D\rvert$, but remains smooth for finite $\beta$.

\subsection{Prime-Regularized Loss}

Given a dataset $\mathcal{D}$, layers $\ell=1,\dots,L$, and bonds $\mathcal{B}$, we define the prime-aware regularization term
\[
\mathcal{L}_{\mathrm{prime}}(\theta)
:= \frac{1}{|\mathcal{D}|}
\sum_{\lvert x \rangle\in \mathcal{D}}
\sum_{\ell=1}^L
\sum_{b\in\mathcal{B}}
\mathrm{dist}_b^{(\ell)}(\theta,x).
\]

The total loss is
\[
\mathcal{L}(\theta)
= \mathcal{L}_{\mathrm{rec}}(\theta)
+ \lambda \mathcal{L}_{\mathrm{prime}}(\theta),
\]
with $\lambda\ge 0$ controlling the strength of rank-alignment to the prime-structured bond menus.

In this weak form, primes influence the training implicitly through the choice of $\mathcal{D}_b$: the model is biased toward effective ranks that match prime-structured capacities.

\section{Prime-Aware Regularization: Strong Form (Exponent Surrogates)}

To more directly encode prime structure in the gradient flow, we can introduce continuous approximations to prime exponent vectors.

\subsection{Log-Based Exponent Surrogates}

Fix a bond $b$ with prime set $\mathcal{P}_b = \{p_{b,1},\dots,p_{b,M_b}\}$. We want to approximate exponents $\tilde{k}_b^{(\ell)}(p_{b,j})$ such that
\[
\log \tilde{r}_b^{(\ell)}
\approx \sum_{j=1}^{M_b} \tilde{k}_b^{(\ell)}(p_{b,j}) \log p_{b,j}.
\]

One can define $\tilde{\mathbf{k}}_b^{(\ell)} = (\tilde{k}_b^{(\ell)}(p_{b,1}),\dots,\tilde{k}_b^{(\ell)}(p_{b,M_b}))$ as the solution of a regularized least-squares problem
\[
\tilde{\mathbf{k}}_b^{(\ell)}
= \arg\min_{\mathbf{k}\in\mathbb{R}^{M_b}}
\left(
\log \tilde{r}_b^{(\ell)} - \sum_{j=1}^{M_b} k_j \log p_{b,j}
\right)^2
+ \mu \sum_{j=1}^{M_b} k_j^2,
\]
which has a closed-form solution in terms of the vector of $\log p_{b,j}$ and the scalar $\log\tilde{r}_b^{(\ell)}$. ($\mu>0$ is a regularization coefficient.)

Because $\tilde{r}_b^{(\ell)}$ is differentiable in $\theta$, this surrogate exponent vector $\tilde{\mathbf{k}}_b^{(\ell)}$ is also differentiable.

\subsection{Exponent-Level Prime Regularization}

Let $\mathbf{k}_b^{\mathrm{target}}$ be a target exponent pattern (e.g.\ favoring powers of $2$ only, or a specific blend of primes). We define
\[
\mathcal{L}_{\mathrm{prime-exp}}(\theta)
= \frac{1}{|\mathcal{D}|}
\sum_{\lvert x \rangle\in \mathcal{D}}
\sum_{\ell=1}^L
\sum_{b\in\mathcal{B}}
\sum_{j=1}^{M_b}
w_{b,j}\,
\left|\tilde{k}_b^{(\ell)}(p_{b,j};\theta,x)
- k_{b,j}^{\mathrm{target}}\right|,
\]
where $w_{b,j}\ge 0$ are weights emphasizing particular primes.

The total loss can then be
\[
\mathcal{L}(\theta)
= \mathcal{L}_{\mathrm{rec}}(\theta)
+ \lambda_1 \mathcal{L}_{\mathrm{prime}}(\theta)
+ \lambda_2 \mathcal{L}_{\mathrm{prime-exp}}(\theta),
\]
expressing both rank-alignment (to discrete capacities) and exponent-alignment (to prime patterns).

In this strong form, the training process explicitly steers not just toward certain rank values, but toward entanglement budgets whose approximate prime-exponent representation matches a pre-specified multiplicity profile.

\section{Prime-Canonical Architecture Selection}

Rather than attempt to define a prime-canonical \emph{TN representation} for a fixed state (which is ill-posed under gauge redundancy and minimality), we define a prime-canonicality concept at the architecture level.

\subsection{Prime-Complexity Functional}

For each bond $b$ and layer $\ell$, let $D_b^{\max,(\ell)}$ denote the maximum allowed bond dimension in $\mathcal{D}_b^{(\ell)}$, with prime exponents $k_{b,j}^{\max,(\ell)}$:
\[
D_b^{\max,(\ell)} = \prod_{j} p_{b,j}^{k_{b,j}^{\max,(\ell)}}.
\]

Define a global \emph{prime-complexity functional}:
\[
C(\{\mathcal{D}_b^{(\ell)}\})
= \sum_{\ell} \sum_{b}
\sum_{j} w_{b,j} k_{b,j}^{\max,(\ell)},
\]
where $w_{b,j}\ge 0$ weight the cost of prime $p_{b,j}$ in bond $b$.

\subsection{Model Selection Criterion}

Given a family of candidate architectures $\{\mathcal{A}_\alpha\}$ (same TN topology, different allowed bond sets $\mathcal{D}_b^{(\ell)}$), we can define a model selection criterion:
\[
\alpha^\star = \arg\min_{\alpha}
\left[
\mathcal{L}_{\mathrm{rec}}^{(\alpha)}(\theta^{(\alpha)\star})
+ \lambda_C C(\mathcal{A}_\alpha)
\right],
\]
where $\theta^{(\alpha)\star}$ is the trained parameter set for architecture $\mathcal{A}_\alpha$. This selects architectures that combine good reconstruction performance with low prime-complexity.

In this sense, a \emph{prime-canonical} architecture is one that lies near the Pareto frontier of reconstruction performance vs.\ prime-structured complexity. This shifts prime-canonicality from a per-state representation notion to a global architectural property.

\section{Illustrative Minimal Example (Conceptual)}

We briefly sketch a minimal 1D example, suitable for an actual numerical test.

\subsection{Example Setup}

\begin{itemize}[leftmargin=1.5em]
    \item Data: a collection of $4$-site binary configurations or low-dimensional quantum states.

    \item Encoder: an MPS-based encoder with three internal bonds $b_1, b_2, b_3$.

    \item Architectures:
    \begin{itemize}
        \item Architecture A (powers of $2$):
        \[
        \mathcal{D}_{b_i} = \{1,2,4\},\quad i=1,2,3.
        \]
        \item Architecture B (introducing prime $3$):
        \[
        \mathcal{D}_{b_i} = \{1,3,6\},\quad i=1,2,3.
        \]
    \end{itemize}
    \item Regularizer: $\mathcal{L}_{\mathrm{prime}}$ as defined above, possibly with \(\lambda > 0\).
\end{itemize}

\subsection{Questions to Probe Experimentally}

\begin{itemize}[leftmargin=1.5em]
    \item Does Architecture B systematically use effective ranks closer to $\{3,6\}$ than Architecture A uses ranks close to $\{2,4\}$?

    \item For fixed parameter budget and similar reconstruction loss, do the two architectures yield qualitatively different latent structures (e.g.\ different clustering in the space of encoded representations)?

    \item Can the strong exponent regularization $\mathcal{L}_{\mathrm{prime-exp}}$ be used to encourage, e.g., predominantly powers-of-two ranks vs.\ mixed primes, and does that qualitatively change training dynamics or interpretability?
\end{itemize}

Actual experiments would require a concrete TN library (ITensor, TeNPy, TensorNetwork, etc.) and suitable hyperparameter choices for $\alpha, \tau, \beta, \lambda$.

\section{Code Snippet (Pseudo-Code / Python-Style Skeleton)}

Below is a minimal pseudo-code sketch in Python style, assuming access to a TN library and an autodiff framework (e.g.\ JAX, PyTorch, or TensorFlow).

\subsection{Rank Surrogate and Prime Regularizer}

\begin{verbatim}
# Pseudo-code; not tied to a specific TN library

def smooth_rank(singular_values, alpha=20.0, tau=0.1):
    # singular_values: 1D array of sigma_j >= 0
    sigma_max = singular_values.max()
    u = singular_values / (sigma_max + 1e-12)
    s = 1.0 / (1.0 + np.exp(-alpha * (u - tau)))
    return s.sum()  # surrogate rank

def soft_dist_to_allowed(r_tilde, allowed_dims, beta=10.0):
    # r_tilde: scalar surrogate rank
    # allowed_dims: list of integers D in D_b
    diffs = np.array([abs(r_tilde - D) for D in allowed_dims])
    weights = np.exp(-beta * diffs)
    weights /= weights.sum()
    return (weights * diffs).sum()

def prime_regularizer(encoder, data, allowed_dims_per_bond, 
                      alpha=20.0, tau=0.1, beta=10.0):
    # encoder: TN encoder object with layered structure
    # data: list of input states |x> (TN or explicit)
    # allowed_dims_per_bond: dict mapping bond_id -> list of integers
    total_reg = 0.0
    count = 0

    for x in data:
        psi = x
        # For each layer in encoder
        for layer in encoder.layers:
            psi = layer.apply(psi)
            psi = canonicalize_mps(psi)  # or TTN canonicalization

            for bond_id in psi.internal_bonds():
                sigmas = psi.singular_values(bond_id)
                r_tilde = smooth_rank(sigmas, alpha=alpha, tau=tau)
                D_allowed = allowed_dims_per_bond[bond_id]
                dist = soft_dist_to_allowed(r_tilde, D_allowed, beta=beta)

                total_reg += dist
                count += 1

    return total_reg / max(count, 1)
\end{verbatim}

The overall loss would then be something like:
\begin{verbatim}
loss = reconstruction_loss(encoder, decoder, data) \
       + lambda_prime * prime_regularizer(encoder, data, allowed_dims_per_bond)
\end{verbatim}

which can be optimized by a standard gradient-based optimizer.

\section{Prior Art and Defensive Publication Positioning}

To position this as prior art, we note the following elements:

\begin{itemize}[leftmargin=1.5em]
    \item \textbf{Tensor-network autoencoders and TN-based ML.} Prior works have developed TN-based autoencoders, quantum-inspired TN autoencoders, and TN frontends in quantum-enhanced or classical ML contexts. These typically treat bond dimensions as hyperparameters or adapt them via generic rank truncation, but do not impose an explicitly prime-structured design nor a prime-aware regularizer as described here.

    \item \textbf{Entanglement regularization and rank control.} There is a growing literature on entanglement regularization and low-rank constraints in tensor networks and neural networks. These works often use entanglement entropy or rank penalties, but not prime-indexed design of \emph{allowed} ranks and prime-exponent-aware regularizers.

    \item \textbf{Prime factorization in TN.} Recent research has begun to explore connections between prime factorization and TNs (e.g.\ for factorization algorithms or structural analysis). However, this is distinct from the present proposal, which uses prime factorization not as a target problem but as a \emph{structural bias} on TN architectures and their learning dynamics.
\end{itemize}

The present document establishes the following distinctive combination as prior art:
\begin{enumerate}[leftmargin=1.5em]
    \item Constraining bond dimensions of a TN autoencoder to discrete sets with specific prime factorizations.

    \item Introducing differentiable rank surrogates based on singular values of canonicalized TN states, and aligning these surrogates to the discrete prime-structured sets via a soft distance regularizer.

    \item Optionally, defining log-based surrogate prime exponent vectors and regularizing them toward target exponent patterns, thereby directly encoding prime-indexed multiplicity constraints into the gradient flow.

    \item Using a prime-exponent-based complexity functional over maximal allowed bond dimensions as a criterion for architecture selection, thereby defining ``prime-canonical'' architectures.

    \item Interpreting multiplicity moves not as runtime topology changes, but as design relationships between allowed bond dimension sets across layers (multiplicity schedules).
\end{enumerate}

\section{Concluding Remarks and Future Extensions}

The NN-AUTOTENSCODER framework developed here is designed to be:

\begin{itemize}[leftmargin=1.5em]
    \item \emph{Mathematically consistent:} all quantities are well-defined; gradients are obtained via standard autodiff; prime structure enters via discrete design and smooth surrogates.

    \item \emph{Computationally implementable:} canonicalization, singular-value extraction, and rank surrogates are standard TN operations; the proposed regularizers are simple to code and differentiate.

    \item \emph{Extendable:} the same ideas can, in principle, be ported to hybrid settings where a TN encoder is mapped to a parameterized quantum circuit (e.g.\ via MPS-to-circuit constructions), with prime-aware architecture constraints analogous at the circuit level. Such extensions would need separate care in gradient estimation (e.g.\ parameter-shift rules), but the prime-structured design principles would carry over.
\end{itemize}

Future directions include:
\begin{itemize}[leftmargin=1.5em]
    \item Empirical evaluation of prime-structured architectures on synthetic and real datasets, including comparisons against standard TN-AEs.

    \item Investigation of whether specific prime patterns improve interpretability, compression, or hardware alignment (e.g.\ matching tensor-core sizes or qubit layouts).

    \item Theoretical analysis of how prime-structured capacity constraints influence expressivity and generalization in TN-based models.

    \item Integration with broader multiplicity theory and state-space models of cognitive or physical systems, where primes index recursively stable multiplicity channels.
\end{itemize}

\appendix
\section*{Mathematical Appendix}
\addcontentsline{toc}{section}{Mathematical Appendix}

\section{Preliminaries and Notation}

We collect here the basic definitions and notational conventions used in the proofs.

\subsection{Tensor-Network Representations and Canonical Forms}

Let $\mathcal{H} = (\mathbb{C}^d)^{\otimes N}$ with $d \ge 2$, $N \in \mathbb{N}$. A pure state $\lvert \psi \rangle \in \mathcal{H}$ is represented as a matrix product state (MPS) with open boundary conditions if there exist tensors
\[
A^{[i]} \in \mathbb{C}^{d \times D_{i-1} \times D_i}, \quad i = 1,\dots,N,
\]
with $D_0 = D_N = 1$ and bond dimensions $D_i \in \mathbb{N}$, such that
\[
\lvert \psi \rangle
= \sum_{s_1,\dots,s_N} 
\left( \sum_{\alpha_1,\dots,\alpha_{N-1}}
A^{[1]}_{s_1,\alpha_1}
A^{[2]}_{s_2,\alpha_1,\alpha_2}
\cdots
A^{[N]}_{s_N,\alpha_{N-1}} \right)
\lvert s_1,\dots,s_N \rangle.
\]

We use the following standard facts (see e.g.\ the literature on canonical forms of MPS):
\begin{itemize}
    \item Any MPS can be brought into a \emph{mixed canonical form} with respect to a given bond $b$ via a sequence of QR and SVD operations.
    \item In this form, across bond $b$, one has a Schmidt decomposition
    \[
    \lvert \psi \rangle
    = \sum_{j=1}^{R_b} \sigma_{b,j} \, \lvert L_{b,j} \rangle \otimes \lvert R_{b,j} \rangle,
    \]
    with singular values $\sigma_{b,1} \ge \cdots \ge \sigma_{b,R_b} > 0$, and $R_b \le \min(D_{b-1}, D_b)$.
    \item The vector $(\sigma_{b,1},\dots,\sigma_{b,R_b})$ is uniquely defined up to ordering.
\end{itemize}

We denote by $\Sigma_b = \mathrm{diag}(\sigma_{b,1},\dots,\sigma_{b,R_b})$ the diagonal matrix of singular values at bond $b$.

\subsection{Operator Norms}

Given a linear operator $T: \mathcal{H} \to \mathcal{H}$, we write $\|T\|_{2}$ for the induced operator norm with respect to the Hilbert space norm,
\[
\|T\|_{2} 
:= \sup_{\|x\|=1} \|T x\|.
\]
If $T$ is represented (in some basis) as a matrix, this coincides with the largest singular value of that matrix.

For two bounded operators $A,B$, we recall the submultiplicative property
\[
\|AB\|_2 \le \|A\|_2 \|B\|_2.
\]

\subsection{Rank Surrogate Function}

For a fixed bond $b$ and layer $\ell$, we defined the \emph{relative} singular values
\[
u_{b,j}^{(\ell)}(\theta,x)
:= \frac{\sigma_{b,j}^{(\ell)}(\theta,x)}{\sigma_{b,\max}^{(\ell)}(\theta,x) + \delta},
\]
where $\sigma_{b,\max}^{(\ell)} := \max_j \sigma_{b,j}^{(\ell)}$ and $\delta > 0$ is a regularization constant. The smooth indicator is
\[
s(u) := \frac{1}{1 + e^{-\alpha (u - \tau)}}, \quad u\in[0,1],
\]
with $\alpha>0$ (steepness) and $\tau\in(0,1)$ (relative threshold). The \emph{rank surrogate} is then
\[
\tilde{r}_b^{(\ell)}(\theta,x)
:= \sum_j s\big(u_{b,j}^{(\ell)}(\theta,x)\big).
\]


\section{Lipschitz Bounds for the Rank Surrogate}

We prove that the rank surrogate is Lipschitz continuous with respect to the singular values, and derive an explicit bound.

\subsection{Lipschitz Continuity of $s(u)$}

\begin{lemma}[Lipschitz bound for $s$]
\label{lem:s-lipschitz}
The function $s(u) = (1 + e^{-\alpha(u-\tau)})^{-1}$ on $\mathbb{R}$ satisfies
\[
\big|s(u) - s(v)\big| \le L_s |u-v|, \quad \forall u,v\in \mathbb{R},
\]
with Lipschitz constant $L_s = \alpha/4$.
\end{lemma}

\begin{proof}
We have
\[
s'(u) = \frac{\alpha e^{-\alpha(u-\tau)}}{\big(1 + e^{-\alpha(u-\tau)}\big)^2}.
\]
Set $z = e^{-\alpha(u-\tau)}$. Then
\[
s'(u) = \frac{\alpha z}{(1+z)^2}.
\]
For $z\ge 0$,
\[
\frac{z}{(1+z)^2} \le \frac{1}{4},
\]
since the maximum of $z/(1+z)^2$ on $[0,\infty)$ occurs at $z=1$ and equals $1/4$. Hence
\[
|s'(u)| \le \frac{\alpha}{4} = L_s, \quad \forall u\in\mathbb{R}.
\]
By the mean value theorem,
\[
|s(u) - s(v)| \le \sup_{w\in[u,v]} |s'(w)|\,|u-v|
\le L_s |u-v|.
\]
\end{proof}

\subsection{Lipschitz Bound for $\tilde{r}_b^{(\ell)}$ in Terms of $u_{b,j}^{(\ell)}$}

\begin{lemma}[Rank surrogate Lipschitz in $u$]
\label{lem:rank-lipschitz-u}
Fix a bond $b$, layer $\ell$, input $x$, and parameters $\theta$. Let
\[
\tilde{r}_b^{(\ell)} = \sum_{j=1}^{R_b} s(u_j),
\]
with $u_j \in [0,1]$. Then for two sets $\{u_j\}$ and $\{v_j\}$ of the same size $R_b$,
\[
\big|\tilde{r}_b^{(\ell)}(\{u_j\}) - \tilde{r}_b^{(\ell)}(\{v_j\})\big|
\le L_s \sum_{j=1}^{R_b} |u_j - v_j|.
\]
\end{lemma}

\begin{proof}
By Lemma~\ref{lem:s-lipschitz},
\[
\big|s(u_j) - s(v_j)\big| \le L_s |u_j - v_j|.
\]
Thus
\[
\big|\tilde{r}_b^{(\ell)}(\{u_j\}) - \tilde{r}_b^{(\ell)}(\{v_j\})\big|
= \left|\sum_{j=1}^{R_b} s(u_j) - \sum_{j=1}^{R_b} s(v_j)\right|
\le \sum_{j=1}^{R_b} |s(u_j) - s(v_j)|
\le L_s \sum_{j=1}^{R_b} |u_j - v_j|.
\]
\end{proof}

\subsection{Dependence on Singular Values}

We now examine the dependence of $u_j = \sigma_j / (\sigma_{\max} + \delta)$ on the singular values $\{\sigma_j\}$.

\begin{lemma}[Lipschitz bound on $u_j$ in terms of $\sigma$]
\label{lem:u-lipschitz}
Let $\bm{\sigma} = (\sigma_1,\dots,\sigma_{R_b})$ and $\bm{\sigma}' = (\sigma_1',\dots,\sigma_{R_b}')$ be two singular-value vectors with nonnegative entries. Let
\[
u_j = \frac{\sigma_j}{\sigma_{\max} + \delta}, \quad
u_j' = \frac{\sigma_j'}{\sigma_{\max}' + \delta},
\]
where $\sigma_{\max} = \max_j \sigma_j$ and $\sigma_{\max}' = \max_j \sigma_j'$, and $\delta > 0$ is fixed. Then there exists a constant $L_u$ (depending on $\delta$ and $R_b$) such that
\[
\sum_{j=1}^{R_b} |u_j - u_j'|
\le L_u \|\bm{\sigma} - \bm{\sigma}'\|_1.
\]
In particular, we can choose
\[
L_u \le \frac{2}{\delta} + \frac{R_b}{\delta^2}\|\bm{\sigma}\|_\infty + \frac{R_b}{\delta^2}\|\bm{\sigma}'\|_\infty.
\]
\end{lemma}

\begin{proof}
We estimate
\[
|u_j - u_j'|
= \left|\frac{\sigma_j}{\sigma_{\max} + \delta} - \frac{\sigma_j'}{\sigma_{\max}' + \delta}\right|
\le \left|\frac{\sigma_j - \sigma_j'}{\sigma_{\max} + \delta}\right|
+ \left|\sigma_j'\right|\left|\frac{1}{\sigma_{\max} + \delta} - \frac{1}{\sigma_{\max}' + \delta}\right|.
\]
First term:
\[
\left|\frac{\sigma_j - \sigma_j'}{\sigma_{\max} + \delta}\right| 
\le \frac{1}{\delta}|\sigma_j - \sigma_j'|,
\]
since $\sigma_{\max} + \delta \ge \delta$.

Second term:
\[
\left|\frac{1}{\sigma_{\max} + \delta} - \frac{1}{\sigma_{\max}' + \delta}\right|
= \frac{|\sigma_{\max}' - \sigma_{\max}|}{(\sigma_{\max} + \delta)(\sigma_{\max}' + \delta)}
\le \frac{1}{\delta^2}|\sigma_{\max} - \sigma_{\max}'|.
\]
Also $|\sigma_{\max} - \sigma_{\max}'|\le \|\bm{\sigma} - \bm{\sigma}'\|_\infty \le \|\bm{\sigma} - \bm{\sigma}'\|_1$. Hence
\[
|u_j - u_j'|
\le \frac{1}{\delta}|\sigma_j - \sigma_j'|
+ \frac{|\sigma_j'|}{\delta^2}\|\bm{\sigma} - \bm{\sigma}'\|_1.
\]
Summing over $j$,
\[
\sum_{j=1}^{R_b} |u_j - u_j'|
\le \frac{1}{\delta}\|\bm{\sigma} - \bm{\sigma}'\|_1
+ \frac{1}{\delta^2}\left(\sum_{j=1}^{R_b}|\sigma_j'|\right)\|\bm{\sigma} - \bm{\sigma}'\|_1.
\]
Using $\sum_j|\sigma_j'|\le R_b \|\bm{\sigma}'\|_\infty$ and an analogous bound involving $\bm{\sigma}$ (if desired symmetrically), we obtain the stated expression for $L_u$.
\end{proof}

\subsection{Combined Lipschitz Bound}

\begin{proposition}[Lipschitz continuity of $\tilde{r}_b^{(\ell)}$ in $\bm{\sigma}$]
\label{prop:rank-lipschitz-sigma}
Under the conditions of Lemmas~\ref{lem:rank-lipschitz-u} and \ref{lem:u-lipschitz}, we have
\[
\big|\tilde{r}_b^{(\ell)}(\bm{\sigma}) - \tilde{r}_b^{(\ell)}(\bm{\sigma}')\big|
\le L_{\mathrm{rank}} \|\bm{\sigma} - \bm{\sigma}'\|_1,
\]
with
\[
L_{\mathrm{rank}} := L_s L_u
= \frac{\alpha}{4} L_u,
\]
where $L_u$ can be taken as in Lemma~\ref{lem:u-lipschitz}.
\end{proposition}

\begin{proof}
Combine Lemma~\ref{lem:rank-lipschitz-u} with Lemma~\ref{lem:u-lipschitz}:
\[
\big|\tilde{r}_b^{(\ell)}(\bm{\sigma}) - \tilde{r}_b^{(\ell)}(\bm{\sigma}')\big|
\le L_s \sum_j |u_j - u_j'|
\le L_s L_u \|\bm{\sigma} - \bm{\sigma}'\|_1.
\]
\end{proof}


\section{Operator-Norm Bounds for Encoder Layers}

We briefly connect changes in singular values to operator perturbations of encoder layers. The goal is to show that small changes in the TN parameters (under mild assumptions) induce small changes in $\tilde{r}_b^{(\ell)}$.

\subsection{Single-Layer Perturbation}

Let $E^{(\ell)}_\theta:\mathcal{H}\to\mathcal{H}$ denote the linear map for encoder layer $\ell$ (for simplicity, consider $\mathcal{H}$ of fixed dimension for an intermediate representation). Suppose we have two parameter settings, $\theta$ and $\theta'$, yielding operators $E^{(\ell)}_\theta$ and $E^{(\ell)}_{\theta'}$.

\begin{lemma}[Singular value perturbation bound]
\label{lem:sv-perturbation}
Fix an input $\lvert \psi \rangle$ with $\|\lvert \psi \rangle\|=1$. Let $\lvert \phi \rangle = E^{(\ell)}_\theta \lvert \psi \rangle$ and $\lvert \phi' \rangle = E^{(\ell)}_{\theta'} \lvert \psi \rangle$. Let $\bm{\sigma}$ and $\bm{\sigma}'$ be the singular values across bond $b$ in mixed canonical form for $\lvert \phi \rangle$ and $\lvert \phi' \rangle$, respectively. Then
\[
\|\bm{\sigma} - \bm{\sigma}'\|_2 
\le \| \lvert \phi \rangle \langle \phi \rvert - \lvert \phi' \rangle \langle \phi' \rvert \|_2
\le 2 \|\lvert \phi - \phi' \rangle\|.
\]
In particular,
\[
\|\bm{\sigma} - \bm{\sigma}'\|_1 
\le \sqrt{R_b}\,\|\bm{\sigma} - \bm{\sigma}'\|_2
\le 2\sqrt{R_b} \|\lvert \phi - \phi' \rangle\|.
\]
\end{lemma}

\begin{proof}
The Schmidt coefficients $\sigma_j$ across a given bipartition correspond to the singular values of the reshaped vector (or equivalently, the eigenvalues of the reduced density matrix). One can bound changes in eigenvalues of the reduced density matrix in terms of operator-norm differences; a simple bound follows from Weyl's inequality or general perturbation bounds. For pure states,
\[
\| \lvert \phi \rangle \langle \phi \rvert - \lvert \phi' \rangle \langle \phi' \rvert \|_2
\le 2 \|\lvert \phi - \phi' \rangle\|,
\]
which follows from direct expansion and the fact that $\|\lvert \phi \rangle \langle \phi \rvert\|_2 = 1$ for normalized vectors. This operator-norm bound then controls the perturbation of the reduced density matrix and hence the Schmidt coefficients; details can be derived from standard matrix perturbation theory.

The inequality $\|\bm{\sigma} - \bm{\sigma}'\|_1 \le \sqrt{R_b}\,\|\bm{\sigma} - \bm{\sigma}'\|_2$ is the Cauchy--Schwarz inequality in $\mathbb{R}^{R_b}$.
\end{proof}

\subsection{Propagating Parameter Perturbations}

Assume the encoder layers are composed:
\[
E_\theta = E^{(L)}_\theta \circ \cdots \circ E^{(1)}_\theta.
\]
For a fixed input $\lvert x \rangle$, denote by $\lvert \psi^{(\ell)}_\theta(x)\rangle$ the state after the $\ell$-th layer and canonicalization (we absorb canonicalization into the definition of $E^{(\ell)}_\theta$ as a map on canonicalized representations).

Under suitable boundedness of each layer (e.g.\ $\|E^{(\ell)}_\theta\|_2 \le C_\ell$), we can bound:

\begin{lemma}[Stability of intermediate states]
\label{lem:state-stability}
Assume $\|E^{(\ell)}_\theta\|_2\le C_\ell$ and $\|E^{(\ell)}_{\theta'}\|_2\le C_\ell$ uniformly in $\theta,\theta'$ for all $\ell$. Let
\[
\Delta^{(\ell)}(x)
:= \|\lvert \psi^{(\ell)}_\theta(x)\rangle - \lvert \psi^{(\ell)}_{\theta'}(x)\rangle\|.
\]
Then
\[
\Delta^{(\ell)}(x) \le 
\left(\prod_{k=1}^{\ell-1} C_k\right)
\|E^{(\ell)}_\theta - E^{(\ell)}_{\theta'}\|_2 \,\|\lvert \psi^{(\ell-1)}_{\theta'}(x)\rangle\|
+ \sum_{m=1}^{\ell-1} 
\left(\prod_{k=m+1}^{\ell} C_k\right)
\|E^{(m)}_\theta - E^{(m)}_{\theta'}\|_2.
\]
In particular, if each $\|E^{(\ell)}_\theta - E^{(\ell)}_{\theta'}\|_2$ is small, then so is $\Delta^{(\ell)}(x)$.
\end{lemma}

\begin{proof}
We proceed by induction, using telescoping sums and the triangle inequality, together with the submultiplicativity $\|AB\|_2 \le \|A\|_2 \|B\|_2$ and the uniform bounds $\|E^{(\ell)}_{\theta}\|_2,\|E^{(\ell)}_{\theta'}\|_2\le C_\ell$. Details follow the standard stability analysis of compositions of bounded linear maps.
\end{proof}

Combining Lemma~\ref{lem:state-stability} with Lemma~\ref{lem:sv-perturbation} and Proposition~\ref{prop:rank-lipschitz-sigma}, we obtain:

\begin{theorem}[Stability of rank surrogate under parameter perturbations]
\label{thm:rank-stability-parameters}
Under the assumptions of Lemma~\ref{lem:state-stability}, for each bond $b$ and layer $\ell$, each input $\lvert x \rangle$, there exists a constant $K_{b,\ell}(x)$ such that
\[
\big|\tilde{r}_b^{(\ell)}(\theta,x) - \tilde{r}_b^{(\ell)}(\theta',x)\big|
\le K_{b,\ell}(x) \max_{1\le m \le \ell} \|E^{(m)}_\theta - E^{(m)}_{\theta'}\|_2.
\]
In particular, if each layer operator depends continuously (or smoothly) on parameters $\theta$, then $\tilde{r}_b^{(\ell)}(\theta,x)$ is continuous (or smooth) in $\theta$.
\end{theorem}

\begin{proof}
From Lemma~\ref{lem:state-stability}, $\Delta^{(\ell)}(x)$ is bounded by a linear combination of $\|E^{(m)}_\theta - E^{(m)}_{\theta'}\|_2$, $m\le \ell$. Lemma~\ref{lem:sv-perturbation} then gives
\[
\|\bm{\sigma}^{(\ell)} - \bm{\sigma}'^{(\ell)}\|_1
\le 2\sqrt{R_b}\,\Delta^{(\ell)}(x).
\]
Proposition~\ref{prop:rank-lipschitz-sigma} yields
\[
\big|\tilde{r}_b^{(\ell)}(\theta,x) - \tilde{r}_b^{(\ell)}(\theta',x)\big|
\le L_{\mathrm{rank}} \|\bm{\sigma}^{(\ell)} - \bm{\sigma}'^{(\ell)}\|_1
\le 2\sqrt{R_b} L_{\mathrm{rank}} \Delta^{(\ell)}(x),
\]
which is then bounded using Lemma~\ref{lem:state-stability}. The constant $K_{b,\ell}(x)$ collects the corresponding multiplicative factors.
\end{proof}


\section{Bounds for the Prime-Regularizer Term}

We now provide simple bounds for the prime-regularizer $\mathcal{L}_{\mathrm{prime}}(\theta)$ and its sensitivity to parameter changes.

\subsection{Bounding $\mathrm{dist}_b^{(\ell)}$}

Recall
\[
\mathrm{dist}_b^{(\ell)}(\theta,x)
= \sum_{i} w_{b,i}^{(\ell)}(\theta,x)\,\big|\tilde{r}_b^{(\ell)}(\theta,x) - D_{b,i}\big|,
\]
where $w_{b,i}^{(\ell)}$ is a softmax over $-\beta|\tilde{r}_b^{(\ell)} - D_{b,i}|$ and $\mathcal{D}_b = \{D_{b,1},\dots,D_{b,m_b}\}$.

\begin{lemma}[Upper bound on $\mathrm{dist}_b^{(\ell)}$]
\label{lem:dist-upper-bound}
For any $\theta,x$, and any bond $b$, layer $\ell$,
\[
\mathrm{dist}_b^{(\ell)}(\theta,x)
\le \max_{i} \big|\tilde{r}_b^{(\ell)}(\theta,x) - D_{b,i}\big|.
\]
Furthermore, if $\tilde{r}_b^{(\ell)}(\theta,x)$ lies in the interval $[D_{b,\min},D_{b,\max}]$, where
\[
D_{b,\min} := \min_i D_{b,i},\quad D_{b,\max}:=\max_i D_{b,i},
\]
then
\[
\mathrm{dist}_b^{(\ell)}(\theta,x)
\le D_{b,\max} - D_{b,\min}.
\]
\end{lemma}

\begin{proof}
The weights $w_{b,i}^{(\ell)}$ form a probability distribution (nonnegative and sum to 1), so $\mathrm{dist}_b^{(\ell)}(\theta,x)$ is a convex combination of $\big|\tilde{r}_b^{(\ell)} - D_{b,i}\big|$. Hence it is bounded above by their maximum. If $\tilde{r}_b^{(\ell)}\in[D_{b,\min},D_{b,\max}]$, then each $|\tilde{r}_b^{(\ell)} - D_{b,i}| \le D_{b,\max} - D_{b,\min}$, giving the second bound.
\end{proof}

\subsection{Lipschitz Continuity of $\mathrm{dist}_b^{(\ell)}$ in $\tilde{r}_b^{(\ell)}$}

\begin{lemma}[Lipschitz bound for $\mathrm{dist}_b^{(\ell)}$]
\label{lem:dist-lipschitz}
Let $r,r'$ be two scalar values, and define
\[
\mathrm{dist}_b(r) = \sum_i w_i(r)\,|r - D_{b,i}|,
\]
with
\[
w_i(r) = \frac{\exp(-\beta |r - D_{b,i}|)}{\sum_k \exp(-\beta |r - D_{b,k}|)}.
\]
There exists a constant $L_{\mathrm{dist}}$ such that
\[
|\mathrm{dist}_b(r) - \mathrm{dist}_b(r')|
\le L_{\mathrm{dist}} |r - r'|,
\]
where $L_{\mathrm{dist}}$ depends on $\beta$ and the set $\mathcal{D}_b$, but not on $r,r'$.
\end{lemma}

\begin{proof}[Sketch]
Both $w_i(r)$ and $|r - D_{b,i}|$ are Lipschitz functions of $r$, with Lipschitz constants bounded in terms of $\beta$ and $\max_i |D_{b,i}|$. Differentiating $\mathrm{dist}_b(r)$ with respect to $r$ and bounding the derivative uniformly in $r$ yields a finite Lipschitz constant $L_{\mathrm{dist}}$. A detailed derivation follows from applying the product rule and bounding each term using the triangle inequality and the properties of the softmax.
\end{proof}

\subsection{Stability of $\mathcal{L}_{\mathrm{prime}}$}

\begin{theorem}[Stability of $\mathcal{L}_{\mathrm{prime}}$ under parameter perturbations]
\label{thm:Lprime-stability}
Under the assumptions of Theorem~\ref{thm:rank-stability-parameters}, for a fixed dataset $\mathcal{D}$, there exists a constant $K_{\mathrm{prime}}$ such that
\[
|\mathcal{L}_{\mathrm{prime}}(\theta) - \mathcal{L}_{\mathrm{prime}}(\theta')|
\le K_{\mathrm{prime}}
\max_{1\le \ell \le L} \|E^{(\ell)}_\theta - E^{(\ell)}_{\theta'}\|_2.
\]
Thus $\mathcal{L}_{\mathrm{prime}}$ is continuous in $\theta$.
\end{theorem}

\begin{proof}
We have
\[
\mathcal{L}_{\mathrm{prime}}(\theta)
= \frac{1}{|\mathcal{D}|}
\sum_{x\in\mathcal{D}}
\sum_{\ell,b}
\mathrm{dist}_b^{(\ell)}(\tilde{r}_b^{(\ell)}(\theta,x)),
\]
where we now emphasize $\mathrm{dist}_b^{(\ell)}$ as a function of $\tilde{r}_b^{(\ell)}$. By Lemma~\ref{lem:dist-lipschitz},
\[
|\mathrm{dist}_b^{(\ell)}(\tilde{r}_b^{(\ell)}(\theta,x)) - \mathrm{dist}_b^{(\ell)}(\tilde{r}_b^{(\ell)}(\theta',x))|
\le L_{\mathrm{dist}} |\tilde{r}_b^{(\ell)}(\theta,x) - \tilde{r}_b^{(\ell)}(\theta',x)|.
\]
Using Theorem~\ref{thm:rank-stability-parameters}, we have
\[
|\tilde{r}_b^{(\ell)}(\theta,x) - \tilde{r}_b^{(\ell)}(\theta',x)|
\le K_{b,\ell}(x) \max_{m\le \ell} \|E^{(m)}_\theta - E^{(m)}_{\theta'}\|_2.
\]
Substituting and summing over $x,\ell,b$, we obtain
\[
|\mathcal{L}_{\mathrm{prime}}(\theta) - \mathcal{L}_{\mathrm{prime}}(\theta')|
\le \frac{1}{|\mathcal{D}|}
\sum_{x,\ell,b}
L_{\mathrm{dist}} K_{b,\ell}(x)
\max_{m\le\ell} \|E^{(m)}_\theta - E^{(m)}_{\theta'}\|_2.
\]
Define
\[
K_{\mathrm{prime}} :=
\frac{1}{|\mathcal{D}|}
\sum_{x,\ell,b}
L_{\mathrm{dist}} K_{b,\ell}(x),
\]
and note that $\max_{m\le\ell}\|E^{(m)}_\theta - E^{(m)}_{\theta'}\|_2 \le \max_{1\le m\le L}\|E^{(m)}_\theta - E^{(m)}_{\theta'}\|_2$. This yields the stated bound.
\end{proof}


\section{Remarks on Strong Prime-Exponent Regularization}

For the strong form $\mathcal{L}_{\mathrm{prime-exp}}$, where approximate exponent vectors $\tilde{\mathbf{k}}_b^{(\ell)}$ are obtained by a regularized least-squares fit,
\[
\tilde{\mathbf{k}}_b^{(\ell)}
= \arg\min_{\mathbf{k}}
\big(\log \tilde{r}_b^{(\ell)} - \sum_j k_j \log p_{b,j}\big)^2
+ \mu \|\mathbf{k}\|_2^2,
\]
we note:

\begin{itemize}
    \item The mapping $\tilde{r} \mapsto \tilde{\mathbf{k}}$ is smooth wherever $\tilde{r}>0$ and the design matrix of $\log p_{b,j}$ has full rank.

    \item The regularizer $\mathcal{L}_{\mathrm{prime-exp}}$ is a composition of linear functionals and absolute values of these $\tilde{\mathbf{k}}_b^{(\ell)}$, hence Lipschitz in $\tilde{r}_b^{(\ell)}$, and therefore continuous in $\theta$ under the same assumptions used for $\mathcal{L}_{\mathrm{prime}}$.

    \item Detailed operator-norm bounds analogous to Theorem~\ref{thm:Lprime-stability} follow by combining:
    \begin{enumerate}
        \item The Lipschitz bounds for $\tilde{r}_b^{(\ell)}$ (Theorem~\ref{thm:rank-stability-parameters});
        \item The smooth dependence of $\tilde{\mathbf{k}}_b^{(\ell)}$ on $\tilde{r}_b^{(\ell)}$ from the closed-form least-squares solution;
        \item A Lipschitz estimate for $\mathcal{L}_{\mathrm{prime-exp}}$ in terms of $\tilde{\mathbf{k}}_b^{(\ell)}$.
    \end{enumerate}
\end{itemize}

For brevity, we omit the full expansion here; the arguments mirror those of the weak form, with additional linear algebra factors arising from the least-squares solution.

\bigskip

\noindent\textbf{Summary.} The results in this appendix establish that:
\begin{itemize}
    \item The rank surrogate $\tilde{r}_b^{(\ell)}$ is Lipschitz continuous in the singular values of the TN state and thus stable under small changes in the encoder parameters.
    \item The prime-regularizer $\mathcal{L}_{\mathrm{prime}}$ inherits this stability and is continuous (indeed, Lipschitz under suitable uniform bounds).
    \item The strong prime-exponent regularizer $\mathcal{L}_{\mathrm{prime-exp}}$ is likewise continuous, with perturbation bounds following from the same structural properties.
\end{itemize}
These properties ensure that the proposed NN-AUTOTENSCODER objective is well behaved under gradient-based optimization and that small changes in parameters cannot cause pathological jumps in the prime-aware regularization terms.


\nocite{*}
\bibliographystyle{unsrt}  
\bibliography{references}
\end{document}